\chapter{Provenance, Reproducibility, and Authority Boundaries}
\label{app:provenance}

\section{Immutable theorem snapshot}

The last immutable Project Koios Formal revision underlying this uncommitted proof-tranche
candidate is merge commit

\begin{center}
\decl{1078380fb608269a8094c7e86776a1c0d826333a},
\end{center}

with Git tree

\begin{center}
\decl{0e9695296ab1275df87e2b4b0309e3ef8c4d24b2}.
\end{center}

The candidate additionally modifies the geometry and finite-compatibility modules.  Its
integration commit and tree do not yet exist and must replace this predecessor boundary
after a separately authorized checkpoint.  Modules central to the alignment, residual,
and compatibility narrative have the following SHA-256 identities in the candidate
working tree:

\begin{description}
  \item[Admissible-set geometry.]
    \modulepath{ProjectKoios/BandGauge/Geometry/AdmissibleSets.lean},
    SHA-256
    \decl{e830b1b940ced88bc86b169261673337126583ee09e9ad9fb30ff580e9160fd8}.
  \item[Finite compatibility schema.]
    \modulepath{ProjectKoios/BandGauge/Decision/FiniteCompatibility.lean},
    SHA-256
    \decl{45e737324df4718497566d57d72ff9bd575f37d59a5f6768a101034b4c232a0d}.
  \item[Aligned operator differences.]
    \modulepath{ProjectKoios/BandGauge/Bloch/AlignedDifference.lean},
    SHA-256
    \decl{050e2cc7736e583c87f65488059f587b259fe5748960822b7f7eaba2485ffab8}.
  \item[Alignment families.]
    \modulepath{ProjectKoios/BandGauge/Alignment/Family.lean},
    SHA-256
    \decl{27b1c817bf10778a6d9032e6d5c8d20601aa9fa12f3c70ab560cdbf7dc672aab}.
  \item[Alignment-family separation.]
    \modulepath{ProjectKoios/BandGauge/Alignment/FamilySeparation.lean},
    SHA-256
    \decl{8811fffaca7411e525f37e710ad6c5b3ea9ba9064cfc74c3435344e8858caa71}.
  \item[Three-channel residual algebra.]
    \modulepath{ProjectKoios/BandGauge/Hilbert/ResidualChannels.lean},
    SHA-256
    \decl{38d0604da95a88a85ecc818037fedda0fb73ee9a51055bffa0a6891d8a4ee150}.
  \item[Gauge--truncation counterexample.]
    \modulepath{ProjectKoios/BandGauge/Fourier/GaugeTruncationCounterexample.lean},
    SHA-256
    \decl{3b4d9083c04592ff9982cbddf891c8243cc599ad6bc32f5e4a10f708bec6186d}.
  \item[Residual channels and finite tails.]
    \modulepath{ProjectKoios/BandGauge/Fourier/ResidualChannels.lean},
    SHA-256
    \decl{d2da775a99edf7dea2d8955a121d4929cc183be4fcf64a515ab5e002c019a28b}.
\end{description}

The predecessor revision and candidate module hashes make the pre-commit boundary
explicit, but they do not yet constitute one immutable integrated source state.  This
candidate is not publication-ready until an authorized checkpoint supplies a commit and
tree and this paragraph is refreshed.  No statement about a later revision follows
merely because a declaration retains the same informal name.  The commit, tree, hashes,
theorem-manifest entries, and proof dossiers must be reviewed together.

\section{Motivating ICMSEP problem statement}

The central spectral--operator compatibility question is taken from the owner repository
\nolinkurl{https://github.com/eragasa/ksdft2effmass.git}.  The source was inspected
read-only on 2026-10-29 at commit

\begin{center}
\decl{4835e7bf9161729cc8f0537bcc13db083896a6ae},
\end{center}

with tree

\begin{center}
\decl{fd641fb8a4efe7a5ef4d5da9c49f126e1e03e355}.
\end{center}

The checked-out branch was
\decl{work/monograph-current-library-capabilities}.  The owner worktree was broadly
dirty, so no repository-wide clean-state claim is made.  The two inspected source files
below were nevertheless byte-identical to the blobs at the recorded commit:

\begin{description}
  \item[Submitted-format extended abstract.]
    \modulepath{docs/publications/conferences/ICMSEP2026/ICMSEP-extended-abstract_02.docx},
    Git blob \decl{d23a7ae3bb137fa6d31f25ace37ba7e25fdd51a7},
    SHA-256
    \decl{4b1a2302ab47f7401b0531d9a89c29dfb05cebfc93e51eab1b8b2c96c351fa13}.
  \item[Markdown companion source.]
    \modulepath{docs/publications/conferences/ICMSEP2026/ksdft2effmass.ICMEP2026.abstract.md},
    Git blob \decl{dfcf5d0ea84a1f81843bd778bea3c8760b66b85d},
    SHA-256
    \decl{d1cf59403dc085d30f17e2bc46db321d9f25840c6f91a762ef157a59bf450ade}.
\end{description}

The extended abstract asks whether spectral and operator criteria admit one common
tight-binding Hamiltonian within declared tolerances, repeats that test over a nested
model hierarchy, and uses admissible-set separation when no common endpoint exists.
That problem statement controls the narrative hierarchy of this companion.  It does not
control Lean theorem types or proof status, and its prospective first-principles and
validation statements are not promoted here into completed scientific results.

\section{Seven-paper error-channel source snapshot}
\label{sec:seven-paper-source-snapshot}

The programmatic account in
Section~\ref{sec:seven-paper-error-channel-development} and
Appendix~\ref{app:seven-paper-error-channel-crosswalk} was inspected read-only in the
same owner repository at commit

\begin{center}
\decl{542ea2c4f16baaf4f3e905c2d833cc95db105dfe},
\end{center}

with tree

\begin{center}
\decl{ae9390fcda3d9142fab83b615104b55b094f8086}.
\end{center}

The checked-out branch was \decl{work/periodic1d-rank1-control-v3}; the worktree and all
listed source paths were clean at inspection on 2026-10-29.  The source identities are:

\begin{description}
  \item[Program roadmap.]
    \modulepath{docs/publications/conferences/ICMSEP2026/paper-program-roadmap.md};
    blob \decl{4db8f5617e7a557378da54e37fffdf03e29e2161};
    SHA-256 \decl{f0cbf0ee15fbaa5f8df80d0fe2d3ffabf02de7cfa3cb34c230c143d3e794b5f5}.
  \item[Experiment-to-paper allocation.]
    \modulepath{docs/publications/papers/experiment-to-paper-allocation.md};
    blob \decl{28d7fe01f176328cdd52268bd8339afce37d202c};
    SHA-256 \decl{e51954b0c2210617235b3e18d57577fe2d76691c099b3cc9d47ef6774fd9dd39}.
  \item[\decl{Ragasa2026a} v3.]
    \modulepath{docs/publications/papers/Ragasa2026a/v3/manuscript.tex};
    blob \decl{ac4dd332e9fc92fcd24c77875cdabbfc6e7852e9};
    SHA-256 \decl{02c20cf087cc7272ba45bb271d617a081582514c618bcdb6fe1c39d389d88d74}.
  \item[\decl{Ragasa2026b} v2.]
    \modulepath{docs/publications/papers/Ragasa2026b/v2/manuscript.tex};
    blob \decl{6e9dfef21e0cb7c33c56f8e01265510dcb9d88e1};
    SHA-256 \decl{d8d5951d8d7ed75355b5d94f37fcbddc70c1466897a5770dbe9f877e60016daf}.
  \item[\decl{Ragasa2026b} diagnostic definitions.]
    \modulepath{docs/publications/papers/Ragasa2026b/v2/appendix_diagnostics.tex};
    blob \decl{e8b62942faea8962cdc8767b1f201962588f3cfa};
    SHA-256 \decl{262b14b192a5893bd99a71d9457664a2c2d498f16af6e545b48273b26ac9e76d}.
  \item[\decl{Ragasa2026c} v1.]
    \modulepath{docs/publications/papers/Ragasa2026c/v1/manuscript.tex};
    blob \decl{4f1bd997ca3296f64472c9788f2668de4d1767bc};
    SHA-256 \decl{baf453722d96cf5b0360a0b80bfa1580465a1db7fe91167583943db10701baf4}.
  \item[\decl{Ragasa2026d} v1.]
    \modulepath{docs/publications/papers/Ragasa2026d/v1/manuscript.tex};
    blob \decl{248cbb6f32e3d1842cfc6be698a2001d1dbd6ff2};
    SHA-256 \decl{bb4aed2614777870c8607c7e67b43dbb133a51c672ac14cec6acf85004a18363}.
  \item[\decl{Ragasa2026e} v1.]
    \modulepath{docs/publications/papers/Ragasa2026e/v1/manuscript.tex};
    blob \decl{c896ffdb4c5ba52f5dbb8185b0738ceaaab755bd};
    SHA-256 \decl{72b60ef8cd2ba5505edeb4acce31ac004721c8ee5f6cee570729e19e870f8b76}.
  \item[\decl{Ragasa2026f} v1.]
    \modulepath{docs/publications/papers/Ragasa2026f/v1/manuscript.tex};
    blob \decl{af1e156b467a3ac4bf39e8b9e5a86f6896a4160a};
    SHA-256 \decl{a11a78b1a5cba72efb4e56695f4b9f1543369e97da7fcc149bee51f53913839c}.
  \item[\decl{Ragasa2026g} v1.]
    \modulepath{docs/publications/papers/Ragasa2026g/v1/manuscript.tex};
    blob \decl{c2b26e248d30f082276a0a3ca41a2b42d951b9ad};
    SHA-256 \decl{c52cc2f73f32c28a95ea2003004d36668b6cfa4c48d103d013c890af1cf42316}.
\end{description}

These owner sources establish editorial scope, terminology, and lifecycle status for the
seven-paper program.  They do not determine Lean theorem types, prove that a numerical
channel was evaluated correctly, or transfer controlled-model evidence into a material
claim.

\section{Authority hierarchy}

Four Formal-owned surfaces have distinct roles:

\begin{enumerate}
  \item Lean source under \modulepath{ProjectKoios/} owns exact definitions,
        propositions, and proof terms.
  \item \modulepath{docs/theorem-manifest.md} owns stable declaration identifiers and
        repository status.
  \item Files under \modulepath{proofs/} provide statement-level explanations,
        assumptions, provenance, and limitations.
  \item This companion synthesizes those declarations into a continuous mathematical
        argument.
\end{enumerate}

External manuscripts, calculation plans, schemas, and research notes may motivate the
formal program. They do not determine theorem statements or proof status. If prose in an
external source and a Lean type differ, the Lean type controls the formal claim and the
editorial crosswalk must record the discrepancy.

\section{Independent status axes}

The following questions remain logically independent:

\begin{itemize}
  \item Is an exact proposition implemented and proved in Lean?
  \item Which axioms does the compiled declaration report?
  \item Has its human-readable proof dossier been reviewed?
  \item Has an external certificate schema been frozen and independently verified?
  \item Does retained numerical evidence satisfy that schema?
  \item Are the theorem's assumptions applicable to the intended scientific system?
  \item Has novelty, publication, or human acceptance been established?
\end{itemize}

A positive answer on one axis cannot be transferred to another without additional
evidence. In particular, a clean Lean build does not certify runtime arrays, and a
reproducible numerical calculation does not prove an exact theorem unless an explicit
checked-certificate bridge connects its retained data to the theorem's premises.

\section{Dependency and toolchain boundary}

The pinned Formal snapshot uses Lean \decl{4.35.0}. Mathlib is pinned to immutable commit
\decl{56a005855ea17e0b0958f73ec801ae140b0a1679}; the dependency inspection reports nine
resolved revisions as commit identities. The durable dependency record is
\modulepath{lake-manifest.json}. A reproducibility review must inspect that file and must
not replace its immutable revisions with floating branches or tags.

The minimum repository validation boundary consists of:

\begin{enumerate}
  \item a successful Lake build for the exact Lean input tree, toolchain, and manifest;
  \item immutable dependency-pin inspection;
  \item an audit for admissions and repository-owned axioms;
  \item coverage of proved declarations by named axiom reports; and
  \item confirmation that generated build state is ignored rather than committed.
\end{enumerate}

A LaTeX build has a different meaning. It establishes that the exposition is
syntactically renderable and permits visual inspection of the resulting pages. It does
not elaborate Lean, validate a scientific model, or authenticate numerical evidence.

\section{Runtime certificate boundary}

The current runtime checked-certificate count for this companion is zero. No binary64
JSON payload, runtime array, eigensolver output, mesh, or serialized result is silently
reinterpreted as exact Lean input.

A future proposition-specific bridge must provide, at minimum:

\begin{itemize}
  \item a frozen schema that rejects unknown fields and missing mandatory data;
  \item exact values or outward-rounded enclosures with a documented rounding policy;
  \item separate producer and independently reviewed verifier implementations;
  \item retained canonical bytes, hashes, source revision, toolchain identity, and
        provenance;
  \item a precise statement of the Lean proposition established when the checker
        accepts; and
  \item strict separation between verification of the certificate and scientific
        assessment of the modeled system.
\end{itemize}

Canonical serialization can authenticate a byte sequence. It cannot make binary64 data
exact, establish enclosure, prove mesh completeness, or certify physical applicability.

\section{Requirements for editorial extraction}

A shortened publication or external chapter derived from this companion must record:

\begin{enumerate}
  \item the exact Formal repository, commit, tree, relevant module paths, and hashes;
  \item every declaration ID used by the extracted argument;
  \item the explicit assumptions and operand, sign, normalization, and domain
        conventions;
  \item the named axiom-audit result and remaining trusted-base limitations;
  \item which surrounding claims remain prospective, numerical, analytical,
        scientific, or publication dependent; and
  \item confirmation that the destination does not present the formal export as runtime
        verification or material validation.
\end{enumerate}

Editorial extraction may reorganize and shorten the exposition. It must not silently
strengthen a theorem, remove a premise, merge distinct alignment families, treat an
unresolved interval as decisive, or detach a statement from its immutable source
identity.
